\chapter*{Acknowledgements}
\addcontentsline{toc}{chapter}{Acknowledgements}

As a web programmer, I made a hundred mistakes a day. If I made less than that, I knew I was behind.  When it comes to computing, mine has been a path of hammering.  I'm the type of person who will hammer and hammer and hammer away at the problem until I eventually break through.  So, I can't promise you a flawless path to success, but I can suggest that we will learn something from this journey.

The mistakes and omissions which follow are entirely mine.  I would like to thank some people for their help along the way.

If I had not found our local hacking group DC423, then I might not have sustained a professional interest in cybersecurity so successfully.  If I had not attended many \index{B-Sides}B-Sides security conferences around the Southeastern \gls{us}, then I might not have realized how grand and stimulating the community of hackers really is.  If I had not continued to go to the library at the University of Tennessee at Chattanooga, the I would not have enjoyed such a full intellectual life after graduation.  If I had not continued my education into adulthood, then my life would certainly have been different.

I had many influential professors along the way.  I fell into the right crowd of instructors at \index{Chattanooga State Community College}Chattanooga State Community College with Leigh MacGregor, Mrs. Peacock, and host of others whose names I cannot recall.  Doctors Kizza, Dumas, Kandah, and especially Li Yang from the \index{University of Tennessee at Chatttanooga}University of Tennessee at Chattanooga built a base that prepared me for later challenges in Computer Science.  Ken Smith, MFA, was my mentor for writing at UTC.  I have many fond memories of the faculty during my undergraduate years there as an English major.  I studied under heroes of the intellect like Doctors Richards, Noe, Rehyansky, Ware, and many others.  \index{Harvard University Extension School}At Harvard Extension School, I'm sure I pissed off almost every TA and instructor who had the misery of grading my work.  For them, I am grateful.  Many thanks to my  professors at Harvard:  Len Evanchik, Daniel Garrie, David Cass, Derek Brink, Ramesh Nagappan, Joseph Ficara, Minlan Yu, David Arias, Eddie Kohler, David Malan, and Bruce Huang.  And, of course, an extra special thanks to the professor for whom I was her "most challenging student," Heather Hinton.

With our chosen style of citations, we tried to provide references that would withstand common lab use and academic rigor.  We provided in-chapter sections for convenient references to books, man pages, and websites.  We also continued with inline footnoting for chapter endnotes; we especially noted various computer commands and technologies since they are fundamentally derivative works upon which we depend.  This made for a lot of references and some duplicity.  We have cited generously.

We decided to set a previous book in fonts by Erik Spiekermann after seeing ``Helvetica'' by Gary Hustwit (2007). We carried that through here.  InfoOTDisp-Book is copyrighted by Ole Schaefer and Erik Spiekermann; it is published by FSI Fonts und Software GmbH.  Eurostile Extended Regular and Eurostile W01 Regular are copyrighted 2006 and 2010 by URW Design and Development.

For typesetting and formatting, we chose LaTex.  We repeatedly consulted various AI and Internet sources for \gls{coding} the document which generated the paper.  For bibliographic formatting, we frequently used MyBib. \endnote{Mybib. “MyBib Bibliography Generator.” MyBib.Com, 13 July 2018, \url{https://www.mybib.com}. Accessed 12 July 2026.}

With the rise of AI, I did choose to use artificial intelligence as I crafted this work.  When I began, near 2015, as a student, it was not available.  As I graduated from Harvard, we often found ourselves in conflict with AI tooling; it was so convenient and effective that most of the time we had to refrain from using it; we were in an intellectual endeavor that required proof of thought; AI is so capable of providing an answer that there are times when it can be a too-helpful assistant.  In our work here, we chose to limit its tasks to mundane data processing, formatting, and organizing text.  Our work could not have effectively been typeset in LaTeX without it.  However, we want to keep the content as intellectually valid as possible so we have chosen to refrain from using it as a tool for composing prose.  We recognize that the whole work is our responsibility.    

\markchapterendnotes